\chapter{Методология и результаты экспериментального исследования}
\section{Протокол и источники данных}
Серии final\_v1 и expanded\_v2 сохранены как исторические результаты. После их аудита выполнены source-only инженерные исправления, а затем проведена серия tuned\_v4 с новым протоколом. Настройки и модели v4 опубликованы до открытия целевых данных~\cite{v4protocol}.

Новые синтетические выборки относятся к известным механизмам генерации. Новые реальные домены --- заранее зарезервированные штаты ACS для задач Income и PublicCoverage. Эти задачи принадлежат общей survey family. Для Beijing и Appliances в v4 выполнена только source-настройка: прежние target уже открыты и не являются независимым тестом новых решений.

\section{Выбор гиперпараметров и парные контроли}
Гиперпараметры выбираются по выделенным исходным доменам, без целевых оценок. Настройка выполняется один раз на development source каждого контекста; затем параметры переносятся на новые обучающие выборки. Сравниваются расширенно настроенный pipeline и ограниченный source-selected reference, а также выбранная модель и её структурно сопоставимый контроль.

Meta uniform использует те же обученные members. MaxRM сохраняет bootstrap и структуру деревьев. У Era контроль с нулевой силой критерия сохраняет доменные ограничения допустимости разбиений и не равен обычному HistGBDT.

\section{Метрики и статистический анализ}
Для классификации используется разность AUROC; для регрессии --- снижение MSE, делённое на дисперсию полного source. Сначала усредняются целевые домены и model seeds, затем сценарии внутри независимой MC replica.

Основные контрасты проверяются двусторонним точным sign test с единой коррекцией Holm. Нулевые контрасты исключаются из sign test. Маргинальный t-интервал среднего не является одновременным интервалом и относится к иному параметру, чем проверка знаков. Реальные перекрывающиеся resamples анализируются описательно~\cite{v4protocol,v4tests}.

\section{Результаты и обсуждение}
Завершено \FinalJobs{} заданий и \PrimaryTests{} основных тестов; после Holm значимы \SignificantTests{} контрастов~\cite{v4audit}. Избранные контрасты представлены в таблице~\ref{tab:v4}. Полная таблица и фильтры строк сохранены в экспериментальных артефактах~\cite{v4tests}.

\begin{table}[htbp]
\caption{Избранные результаты v4; источник: \cite{v4tests}}
\label{tab:v4}
\centering
\small
\begin{tabularx}{\textwidth}{@{}lXrr@{}}
\toprule
Метод & Контраст и задача & Эффект & Holm $p$ \\
\midrule
IRF & Настройка, классификация & +0.035031 & 0.046875 \\
CRF & Настройка, классификация & +0.035034 & 0.046875 \\
META & Механизм, классификация & +0.002372 & 0.046875 \\
ERA & Механизм, регрессия & +0.014899 & 0.046875 \\
IRF & Механизм, классификация & -0.000431 & 0.046875 \\
CRF & Механизм, классификация & -0.001086 & 0.046875 \\
\bottomrule
\end{tabularx}
\par\smallskip\footnotesize Классификация: разность AUROC; регрессия: снижение MSE / source variance.
\end{table}


Расширение настройки улучшило ряд моделей относительно ограниченного reference. При этом положительный вклад механизма подтверждён не у всех методов. Большой HPO gain не означает, что выигрыш обусловлен DG-регуляризацией, и не гарантирует победы над сильным baseline.

\section{Ограничения исследования}
Синтетические MC replicas не являются выборкой всех табличных задач. Ограниченный поиск не устанавливает глобальный максимум каждого метода; CPU бюджеты различаются. Положительные средние контрасты не означают устойчивости при каждом отдельном сдвиге. Локальные hashes и timestamps не подменяют независимый системный аудит доступа.

\section{Вывод по экспериментальной части}
Экспериментальные данные обосновывают раздельную оценку настройки, механизмов и абсолютного качества. Окончательная редакция главы должна включать полный разбор v1/v2/v4, таблицы реальных задач и анализ каждого вида сдвига; исходные материалы сохранены в репозитории.
